%%%%%%%%%%%%%%%%%%%%%%%%%%%%%%%%%%%%%%%%%%%%%%%%%%%%%%%%%%%%
\section{Randnotizen und TODO-Notizen}%
\index{Randnotizen}%
\index{TODO-Notizen}%
\label{sec:Randnotizen}
%%%%%%%%%%%%%%%%%%%%%%%%%%%%%%%%%%%%%%%%%%%%%%%%%%%%%%%%%%%%
%
Randnotizen 
\floatmarginnote{Ich bin eine überflüssige Randnotiz}%
werden mit dem Kommando \lc{floatmarginnote} gesetzt.
Sie eignen sich zum Beispiel, um wichtige Begriffe oder Aussagen zu verdeutlichen
oder durch eine Kurzzusammenfassung jedes einzelnen Textabschnitts
den roten Faden zu verdeutlichen.
Aufgrund des sehr kleinen Randbereichs im Seitenlayout
können in der endgültigen Fassung des Manuskriptes keine Randnotizen platziert werden.
Um \ggf vorhandene Randnotizen auszublenden,
muss man in der Datei \printfilepathquoted{preambel/AlleSchalter.tex} den Wert des Schalters
\printkeyword{showMarginNotes} auf \printkeyword{false} setzen.

TODOs im Text lassen sich mit Hilfe des Befehls \lc{todo\{<Hinweistext>\}} aus dem Paket \pkg{todonotes} setzen.
\todo{Hier weitere Beispiele hinzufügen!}

Mit dem Befehl \lc{missingfigure\{<Hinweistext>\}} lässt sich auf eine fehlende Grafik hinweisen.
\missingfigure{Hier fehlt eine Grafik!}

Die TODOs lassen sich ausblenden, indem man in der Datei \printfilepathquoted{preambel/AlleSchalter.tex} den Wert des Schalters
\printkeyword{showTODOs} auf \printkeyword{false} setzt.

Sofern TODOs vorhanden und eingeblendet sind, wird am Ende des Manuskriptes automatisch eine \hyperlink{TODO-List}{Übersicht der TODOs} engefügt.